\chapter{2007年福建卷（文）}

\section{单选题}

\begin{problem}
已知全集 \(U = \{1, 2, 3, 4, 5\}\)，且 \(A = \{2, 3, 4\}\)，\(B = \{1, 2\}\)，则 \(A \cap \left( \complement_U B \right)\) 等于（\quad）

\choices
    {\(\{2\}\)}
    {\(\{5\}\)}
    {\(\{3, 4\}\)}
    {\(\{2, 3, 4, 5\}\)}
\end{problem}

\begin{answer}
C
\end{answer}

\begin{solution}
由全集定义，\(\complement_U B=\{3,4,5\}\)，所以
\[
A\cap\left(\complement_U B\right)=\{2,3,4\}\cap\{3,4,5\}=\{3,4\}
\]
因此选 C.
\end{solution}

\begin{problem}
等比数列 \(\{a_n\}\) 中，\(a_4 = 4\)，则 \(a_2 \cdot a_6\) 等于（\quad）

\choices
    {\(4\)}
    {\(8\)}
    {\(16\)}
    {\(32\)}
\end{problem}

\begin{answer}
C
\end{answer}

\begin{solution}
设等比数列的公比为 \(q\). 由 \(a_4=a_2q^2\)，\(a_6=a_4q^2\)，得
\[
a_2a_6=a_2a_4q^2=a_4^2=4^2=16
\]
因此选 C.
\end{solution}

\begin{problem}
\(\sin 15^\circ \cos 75^\circ + \cos 15^\circ \sin 105^\circ\) 等于（\quad）

\choices
    {\(0\)}
    {\(\frac{1}{2}\)}
    {\(\frac{\sqrt{3}}{2}\)}
    {\(1\)}
\end{problem}

\begin{answer}
D
\end{answer}

\begin{solution}
利用诱导公式，\(\cos75^\circ=\sin15^\circ\)，\(\sin105^\circ=\cos15^\circ\). 因而
\[
\sin15^\circ\cos75^\circ+\cos15^\circ\sin105^\circ
=\sin^2 15^\circ+\cos^2 15^\circ=1
\]
因此选 D.
\end{solution}

\begin{problem}
“\(|x| < 2\)”是“\(x^2 - x - 6 < 0\)”的（\quad）

\choices
    {充分而不必要条件}
    {必要而不充分条件}
    {充要条件}
    {既不充分也不必要条件}
\end{problem}

\begin{answer}
A
\end{answer}

\begin{solution}
由 \(|x|<2\) 得 \(-2<x<2\). 又
\[
x^2-x-6=(x-3)(x+2)
\]
所以 \(x^2-x-6<0\) 的解集为 \((-2,3)\). 因此 \(|x|<2\) 能推出 \(x^2-x-6<0\)，但反之不成立，例如 \(x=\frac52\) 满足 \(-2<x<3\)，却不满足 \(|x|<2\). 故前者是后者的充分而不必要条件，选 A.
\end{solution}

\begin{problem}
函数 \(f(x) = \sin \left( 2x + \frac{\pi}{3} \right)\) 的图象（\quad）

\choices
    {关于点 \(\left( \frac{\pi}{3}, 0 \right)\) 对称}
    {关于直线 \(x = \frac{\pi}{4}\) 对称}
    {关于点 \(\left( \frac{\pi}{4}, 0 \right)\) 对称}
    {关于直线 \(x = \frac{\pi}{3}\) 对称}
\end{problem}

\begin{answer}
A
\end{answer}

\begin{solution}
对任意实数 \(u\)，有
\[
f\left(\frac{\pi}{3}+u\right)=\sin\left(\pi+2u\right)=-\sin2u
\]
\[
f\left(\frac{\pi}{3}-u\right)=\sin\left(\pi-2u\right)=\sin2u
\]
两式相加为零，故图象上这两个横坐标相反于 \(\frac{\pi}{3}\) 的点的纵坐标也相反，它们的中点为 \(\left(\frac{\pi}{3},0\right)\). 因此图象关于点 \(\left(\frac{\pi}{3},0\right)\) 对称，选 A.
\end{solution}

\begin{problem}
如图，在正方体 \(ABCD - A_1B_1C_1D_1\) 中，\(E\)，\(F\)，\(G\)，\(H\) 分别为 \(AA_1\)，\(AB\)，\(BB_1\)，\(B_1C_1\) 的中点，则异面直线 \(EF\) 与 \(GH\) 所成的角等于（\quad）

\bitmapfigure[max width=0.35\linewidth]{img/2007/fujian_liberal/q6_fig1.png}

\choices
    {\(45^\circ\)}
    {\(60^\circ\)}
    {\(90^\circ\)}
    {\(120^\circ\)}
\end{problem}

\begin{answer}
B
\end{answer}

\begin{solution}
取正方体棱长为 \(2\)，建立直角坐标系：
\(A=(0,0,0)\)，\(B=(2,0,0)\)，\(C=(2,2,0)\)，\(D=(0,2,0)\)
且 \(A_1=(0,0,2)\)，\(B_1=(2,0,2)\)，\(C_1=(2,2,2)\). 于是
\(E=(0,0,1)\)，\(F=(1,0,0)\)，\(G=(2,0,1)\)，\(H=(2,1,2)\).
取两直线的方向向量
\(\overrightarrow{EF}=(1,0,-1)\)，\(\overrightarrow{GH}=(0,1,1)\).
它们的内积为 \(-1\)，长度均为 \(\sqrt2\). 设异面直线所成的锐角为 \(\theta\)，则
\[
\cos\theta=\frac{|\overrightarrow{EF}\cdot\overrightarrow{GH}|}{|\overrightarrow{EF}|\,|\overrightarrow{GH}|}=\frac{1}{2}
\]
所以 \(\theta=60^\circ\)，选 B.
\end{solution}

\begin{problem}
已知 \(f(x)\) 为 \(\R\) 上的减函数，则满足 \(f\left( \frac{1}{x} \right) > f(1)\) 的实数 \(x\) 的取值范围是（\quad）

\choices
    {\((-\infty, 1)\)}
    {\((1, +\infty)\)}
    {\((-\infty, 0) \cup (0, 1)\)}
    {\((-\infty, 0) \cup (1, +\infty)\)}
\end{problem}

\begin{answer}
D
\end{answer}

\begin{solution}
因 \(f\) 在 \(\R\) 上为减函数，故
\[
f\left(\frac1x\right)>f(1)\Longleftrightarrow \frac1x<1
\]
同时必须有 \(x\neq0\). 当 \(x<0\) 时，不等式 \(\frac1x<1\) 恒成立；当 \(x>0\) 时，等价于 \(x>1\). 因此
\[
x\in(-\infty,0)\cup(1,+\infty)
\]
选 D.
\end{solution}

\begin{problem}
对于向量 \(\bs{a}\)，\(\bs{b}\)，\(\bs{c}\) 和实数 \(\lambda\)，下列命题中真命题的是（\quad）

\choices
    {若 \(\bs{a} \cdot \bs{b} = 0\)，则 \(\bs{a} = 0\) 或 \(\bs{b} = 0\)}
    {若 \(\lambda\bs{a} = 0\)，则 \(\lambda = 0\) 或 \(\bs{a} = 0\)}
    {若 \(\bs{a}^2 = \bs{b}^2\)，则 \(\bs{a} = \bs{b}\) 或 \(\bs{a} = -\bs{b}\)}
    {若 \(\bs{a} \cdot \bs{b} = \bs{a} \cdot \bs{c}\)，则 \(\bs{b} = \bs{c}\)}
\end{problem}

\begin{answer}
B
\end{answer}

\begin{solution}
若 \(\lambda\bs{a}=0\) 且 \(\lambda\neq0\)，则两边同除以非零实数 \(\lambda\)，得到 \(\bs{a}=0\)，所以必有 \(\lambda=0\) 或 \(\bs{a}=0\)，命题 B 正确.

命题 A 不成立，例如取 \(\bs{a}=(1,0)\)，\(\bs{b}=(0,1)\)，则 \(\bs{a}\cdot\bs{b}=0\)，但两向量都不是零向量. 命题 C 仅能说明两向量长度相等，例如上述两个向量长度相等，却既不相等也不互为相反向量. 命题 D 也不成立，例如取 \(\bs{a}=(1,0)\)，\(\bs{b}=(0,1)\)，\(\bs{c}=(0,2)\)，则 \(\bs{a}\cdot\bs{b}=\bs{a}\cdot\bs{c}=0\)，但 \(\bs{b}\neq\bs{c}\). 因此选 B.
\end{solution}

\begin{problem}
已知 \(m\)，\(n\) 为两条不同的直线，\(\alpha\)，\(\beta\) 为两个不同的平面，则下列命题中正确的是（\quad）

\choices
    {\(m \subset \alpha\)，\(n \subset \alpha\)，\(m \parallel \beta\)，\(n \parallel \beta \Rightarrow \alpha \parallel \beta\)}
    {\(\alpha \parallel \beta\)，\(m \subset \alpha\)，\(n \subset \beta \Rightarrow m \parallel n\)}
    {\(m \perp \alpha\)，\(m \perp n \Rightarrow n \parallel \alpha\)}
    {\(n \parallel m\)，\(n \perp \alpha \Rightarrow m \perp \alpha\)}
\end{problem}

\begin{answer}
D
\end{answer}

\begin{solution}
命题 D 中，若 \(n\parallel m\) 且 \(n\perp\alpha\)，则与 \(n\) 平行的直线 \(m\) 也垂直于平面 \(\alpha\)，所以 D 正确.

其余命题均不能成立. 例如在 \(\alpha:z=0\)，\(\beta:y=0\) 时，取 \(m:y=1\)，\(z=0\) 和 \(n:y=2\)，\(z=0\)，则 \(m\)，\(n\subset\alpha\) 且都平行于 \(\beta\)，但 \(\alpha\) 与 \(\beta\) 相交，故 A 不成立. 取 \(\alpha:z=0\)，\(\beta:z=1\)，再取 \(m\) 为 \(\alpha\) 内的 \(x\) 轴、\(n\) 为 \(\beta\) 内的 \(y\) 轴，则 \(m\) 与 \(n\) 不平行，故 B 不成立. 对 C，取 \(m\) 为 \(z\) 轴，\(\alpha:z=0\)，取 \(n\) 为 \(\alpha\) 内的 \(x\) 轴，则 \(m\perp\alpha\) 且 \(m\perp n\)，但 \(n\subset\alpha\) 而不是 \(n\parallel\alpha\). 因此选 D.
\end{solution}

\begin{problem}
以双曲线 \(x^2 - y^2 = 2\) 的右焦点为圆心，且与其右准线相切的圆的方程是（\quad）

\choices
    {\(x^2 + y^2 - 4x - 3 = 0\)}
    {\(x^2 + y^2 - 4x + 3 = 0\)}
    {\(x^2 + y^2 + 4x - 5 = 0\)}
    {\(x^2 + y^2 + 4x + 5 = 0\)}
\end{problem}

\begin{answer}
B
\end{answer}

\begin{solution}
双曲线可写为
\[
\frac{x^2}{2}-\frac{y^2}{2}=1
\]
故 \(a^2=2\)，\(b^2=2\)，\(c^2=a^2+b^2=4\)，右焦点为 \((2,0)\)，右准线方程为
\[
x=\frac{a^2}{c}=1
\]
以 \((2,0)\) 为圆心且与直线 \(x=1\) 相切的圆，半径为圆心到准线的距离 \(1\)，所以
\[
(x-2)^2+y^2=1
\]
即 \(x^2+y^2-4x+3=0\). 因此选 B.
\end{solution}

\begin{problem}
已知对任意实数 \(x\)，有 \(f(-x) = -f(x)\)，\(g(-x) = g(x)\)，且 \(x > 0\) 时，\(f'(x) > 0\)，\(g'(x) > 0\)，则 \(x < 0\) 时（\quad）

\choices
    {\(f'(x) > 0\)，\(g'(x) > 0\)}
    {\(f'(x) > 0\)，\(g'(x) < 0\)}
    {\(f'(x) < 0\)，\(g'(x) > 0\)}
    {\(f'(x) < 0\)，\(g'(x) < 0\)}
\end{problem}

\begin{answer}
B
\end{answer}

\begin{solution}
对 \(f(-x)=-f(x)\) 两边求导，得
\[
-f'(-x)=-f'(x)
\]
所以 \(f'(-x)=f'(x)\). 当 \(x<0\) 时 \(-x>0\)，从而 \(f'(x)=f'(-x)>0\).

对 \(g(-x)=g(x)\) 两边求导，得
\[
-g'(-x)=g'(x)
\]
所以 \(g'(x)=-g'(-x)<0\). 因此选 B.
\end{solution}

\begin{problem}
某通讯公司推出一组手机号码，卡号的前七位数字固定，从“××××××× \(0000\)”到“×××××××\(9999\)”共 \(10000\) 个号码. 公司限定：凡卡号的后四位带有数字“\(4\)”或“\(7\)”的一律作为“优惠卡”，则这组号码中“优惠卡”的个数为（\quad）

\choices
    {\(2000\)}
    {\(4096\)}
    {\(5904\)}
    {\(8320\)}
\end{problem}

\begin{answer}
C
\end{answer}

\begin{solution}
后四位每一位都有 \(10\) 种取法，共有 \(10^4=10000\) 个号码. 不含数字 \(4\) 和 \(7\) 时，每一位有 \(8\) 种取法，共有 \(8^4=4096\) 个号码. 因而至少含有一个数字 \(4\) 或 \(7\) 的号码数为
\[
10000-4096=5904
\]
因此选 C.
\end{solution}

\section{填空题}

\begin{problem}
\(\left( x^2 + \frac{1}{x} \right)^6\) 的展开式中常数项是\fillinblank{}. （用数字作答）
\end{problem}

\begin{answer}
15
\end{answer}

\begin{solution}
在展开式中，设从六个因式中选出 \(k\) 个取 \(x^2\)，其余取 \(x^{-1}\)，则对应项的幂为
\[
2k-(6-k)=3k-6
\]
要得到常数项，必须有 \(3k-6=0\)，即 \(k=2\). 因此常数项的系数为
\[
\mathrm C_6^2=15
\]
\end{solution}

\begin{problem}
已知实数 \(x\)，\(y\) 满足 \(\begin{cases}
x + y \geqslant 2,\\
x - y \leqslant 2,\\
0 \leqslant y \leqslant 3.
\end{cases}\)，则 \(2x - y\) 的取值范围是\fillinblank{}.
\end{problem}

\begin{answer}
\(\left[-5, 7\right]\)
\end{answer}

\begin{solution}
由约束条件得
\(2-y\leqslant x\leqslant y+2\)，\(0\leqslant y\leqslant3\).
因此对固定的 \(y\)，有
\[
4-3y=2(2-y)-y\leqslant2x-y\leqslant2(y+2)-y=y+4
\]
当 \(0\leqslant y\leqslant3\) 时，左端的最小值为 \(-5\)，右端的最大值为 \(7\). 取 \(y=3\) 时，\(-1\leqslant x\leqslant5\)，从而 \(2x-y\) 恰好取得区间 \([-5,7]\) 内的所有值. 所以取值范围为 \(\left[-5,7\right]\).
\end{solution}

\begin{problem}
已知长方形 \(ABCD\)，\(AB = 4\)，\(BC = 3\)，则以 \(A\)，\(B\) 为焦点，且过 \(C\)，\(D\) 两点的椭圆的离心率为\fillinblank{}.
\end{problem}

\begin{answer}
\(\frac{1}{2}\)
\end{answer}

\begin{solution}
在长方形中，\(AC=\sqrt{AB^2+BC^2}=\sqrt{4^2+3^2}=5\)，\(BC=3\). 因为 \(A\)，\(B\) 是椭圆的两个焦点且 \(C\) 在椭圆上，所以
\(2a=AC+BC=5+3=8\)，\(a=4\).
又两焦点间距离 \(AB=2c=4\)，故 \(c=2\). 因此离心率
\[
e=\frac ca=\frac24=\frac12
\]
\end{solution}

\begin{problem}
中学数学中存在许多关系，比如“相等关系”“平行关系”等等. 如果集合 \(A\) 中元素之间的一个关系“\(\sim\)”满足以下三个条件：

\begin{circlelist}
\item 自反性：对于任意 \(a \in A\)，都有 \(a \sim a\).
\item 对称性：对于任意 \(a\)，\(b \in A\)，若 \(a \sim b\)，则有 \(b \sim a\).
\item 传递性：对于任意 \(a\)，\(b\)，\(c \in A\)，若 \(a \sim b\)，\(b \sim c\)，则有 \(a \sim c\).
\end{circlelist}

则称“\(\sim\)”是集合 \(A\) 的一个等价关系. 例如：“数的相等”是等价关系，而“直线的平行”不是等价关系（自反性不成立）. 请你再列出两个等价关系：\fillinblank{}.
\end{problem}

\begin{answer}
\(\text{图形的全等、图形的相似}\)
\end{answer}

\begin{solution}
可取“图形的全等关系”和“图形的相似关系”. 对任意图形，它都与自身全等、相似，所以满足自反性；若图形 \(X\) 与图形 \(Y\) 全等（或相似），则交换两者仍成立，所以满足对称性；若图形 \(X\) 与图形 \(Y\) 全等（或相似），且图形 \(Y\) 与图形 \(Z\) 全等（或相似），则将对应的全等变换合成（或将相似比相乘）可得图形 \(X\) 与图形 \(Z\) 全等（或相似），所以满足传递性. 因此这两个关系都是等价关系.
\end{solution}

\section{解答题}

\begin{problem}
在 \(\triangle ABC\) 中，\(\tan A = \frac{1}{4}\)，\(\tan B = \frac{3}{5}\).

\begin{enumerate}
\item 求角 \(C\) 的大小.
\item 若 \(AB\) 边的长为 \(\sqrt{17}\)，求 \(BC\) 边的长.
\end{enumerate}
\end{problem}

\begin{solution}
\begin{enumerate}
\item 因为 \(A\)，\(B\) 是三角形的内角，且 \(\tan A>0\)，\(\tan B>0\)，所以 \(A\)，\(B\) 都是锐角. 因此
\[
\tan(A+B)=\frac{\tan A+\tan B}{1-\tan A\tan B}
 =\frac{\frac14+\frac35}{1-\frac14\cdot\frac35}=1
\]
又 \(0<A+B<\pi\)，且 \(\tan(A+B)>0\)，所以 \(0<A+B<\frac{\pi}{2}\)，从而 \(A+B=\frac{\pi}{4}\). 因此
\[
C=\pi-(A+B)=\frac{3\pi}{4}=135^{\circ}
\]
\item 由 \(\tan A=\frac14\) 且 \(A\) 为锐角，得
\[
\sin A=\frac{\tan A}{\sqrt{1+\tan^2 A}}=\frac{1}{\sqrt{1+4^{2}}}=\frac{1}{\sqrt{17}}
\]
由正弦定理，
\[
\frac{BC}{\sin A}=\frac{AB}{\sin C}
\]
所以
\[
BC=AB\frac{\sin A}{\sin C}
 =\sqrt{17}\cdot\frac{1}{\sqrt{17}}\cdot\frac{1}{\frac{\sqrt{2}}{2}}
 =\sqrt{2}
\]
故 \(BC=\sqrt{2}\).
\end{enumerate}
\end{solution}

\begin{problem}
甲，乙两名跳高运动员一次试跳 \(2\) 米高度成功的概率分别为 \(0.7\)，\(0.6\)，且每次试跳成功与否相互之间没有影响，求：

\begin{enumerate}
\item 甲试跳三次，第三次才成功的概率.
\item 甲，乙两人在第一次试跳中至少有一人成功的概率.
\item 甲，乙各试跳两次，甲比乙的成功次数恰好多一次的概率.
\end{enumerate}
\end{problem}

\begin{solution}
\begin{enumerate}
\item 甲前两次试跳失败、第三次试跳成功，且各次试跳相互独立，所以所求概率为
\[
P=(1-0.7)^2\cdot0.7=0.063
\]
\item 甲、乙两人在第一次试跳中至少有一人成功的对立事件是两人都失败，因此所求概率为
\[
P=1-(1-0.7)(1-0.6)=1-0.3\cdot0.4=0.88
\]
\item 设甲、乙各试跳两次的成功次数分别为 \(X\)，\(Y\). 事件 \(X=Y+1\) 分为互斥的两种情形：
\[
(X=1,Y=0)\quad\text{或}\quad(X=2,Y=1)
\]
因此
\[
P(X=Y+1)=\mathrm C_2^1(0.7)(0.3)\cdot(0.4)^2
 +(0.7)^2\cdot\mathrm C_2^1(0.6)(0.4)
\]
\[
P(X=Y+1)=0.42\cdot0.16+0.49\cdot0.48=0.3024
\]
故所求概率为 \(0.3024\).
\end{enumerate}
\end{solution}

\begin{problem}
如图，正三棱柱 \(ABC - A_1B_1C_1\) 的所有棱长都为 \(2\)，\(D\) 为 \(CC_1\) 中点.

\begin{enumerate}
\item 求证：\(AB_1 \perp\) 平面 \(A_1BD\).
\item 求二面角 \(A-A_1D-B\) 的大小.
\end{enumerate}

\bitmapfigure[max width=0.42\linewidth]{img/2007/fujian_liberal/q19_fig1.png}
\end{problem}

\begin{solution}
\begin{enumerate}
\item 建立空间直角坐标系，取 \(A=(0,0,0)\)，\(B=(2,0,0)\)，\(C=(1,\sqrt{3},0)\). 另取 \(A_1=(0,0,2)\)，\(B_1=(2,0,2)\)，\(C_1=(1,\sqrt{3},2)\). 因为 \(D\) 是 \(CC_1\) 的中点，所以 \(D=(1,\sqrt{3},1)\). 从而
\(\overrightarrow{AB_1}=(2,0,2)\)，\(\overrightarrow{BA_1}=(-2,0,2)\)，\(\overrightarrow{BD}=(-1,\sqrt{3},1)\).
有
\(\overrightarrow{AB_1}\cdot\overrightarrow{BA_1}=-4+4=0\)，\(\overrightarrow{AB_1}\cdot\overrightarrow{BD}=-2+2=0\).
再由
\[
\overrightarrow{BA_1}\times\overrightarrow{BD}=(-2\sqrt{3},0,-2\sqrt{3})
\]
知 \(\overrightarrow{AB_1}\) 平行于平面 \(A_1BD\) 的法向量. 该平面过点 \(B\)，故方程为 \(x+z=2\)；直线 \(AB_1\) 可表示为 \((x,y,z)=(2s,0,2s)\)，在 \(s=\frac12\) 时经过点 \((1,0,1)\)，且该点满足平面方程. 所以直线 \(AB_1\) 与平面 \(A_1BD\) 相交且方向垂直于该平面，即
\[
AB_1\perp\text{平面 }A_1BD
\]
\item 设二面角 \(A-A_1D-B\) 的大小为 \(\theta\). 其两个半平面所在的平面分别为 \(AA_1D\) 和 \(BA_1D\).

平面 \(AA_1D\) 的一个法向量和平面 \(A_1BD\) 的一个法向量分别可取
\[
n_1=\overrightarrow{AA_1}\times\overrightarrow{AD}
 =(-2\sqrt{3},2,0)
\]
\[
n_2=\overrightarrow{A_1B}\times\overrightarrow{A_1D}
 =(2\sqrt{3},0,2\sqrt{3})
\]
因此 \(\lVert n_1\rVert=4\)，\(\lVert n_2\rVert=2\sqrt{6}\)，\(\left|n_1\cdot n_2\right|=12\).
由两平面所成角与其法向量所成锐角相等，得
\[
\cos\theta=\frac{12}{4\cdot2\sqrt{6}}=\frac{\sqrt{6}}{4}
\]
所以
\[
\theta=\arccos\frac{\sqrt{6}}{4}
\]
\end{enumerate}
\end{solution}

\begin{problem}
设函数 \(f(x) = tx^2 + 2t^2x + t - 1\)（\(x \in \R\)，\(t > 0\)）.

\begin{enumerate}
\item 求 \(f(x)\) 的最小值 \(h(t)\).
\item 若 \(h(t) < -2t + m\) 对 \(t \in (0, 2)\) 恒成立，求实数 \(m\) 的取值范围.
\end{enumerate}
\end{problem}

\begin{solution}
\begin{enumerate}
\item 配方得
\[
f(x)=t\left(x+t\right)^2-t^3+t-1
\]
因为 \(t>0\)，所以当 \(x=-t\) 时，\(f(x)\) 取得最小值，故
\[
h(t)=-t^3+t-1
\]
\item 条件
\(h(t)<-2t+m\)，\(\left(t\in(0,2)\right)\)
等价于
\(m>-t^3+3t-1\)，\(\left(t\in(0,2)\right)\).
令 \(\varphi(t)=-t^3+3t-1\)，则
\[
\varphi'(t)=3(1-t^2)
\]
所以 \(\varphi(t)\) 在 \((0,1)\) 上递增，在 \((1,2)\) 上递减，且
\[
\max_{0<t<2}\varphi(t)=\varphi(1)=1
\]
由于 \(t=1\) 属于给定区间且原不等式为严格不等式，必须有 \(m>1\). 反之，当 \(m>1\) 时，对任意 \(t\in(0,2)\) 都有 \(m>\varphi(t)\)，条件成立.

故 \(m\) 的取值范围是 \((1,+\infty)\).
\end{enumerate}
\end{solution}

\begin{problem}
数列 \(\{a_n\}\) 的前 \(n\) 项和为 \(S_n\)，\(a_1 = 1\)，\(a_{n+1} = 2S_n\)（\(n \in \N^*\)）.

\begin{enumerate}
\item 求数列 \(\{a_n\}\) 的通项 \(a_n\).
\item 求数列 \(\{na_n\}\) 的前 \(n\) 项和 \(T_n\).
\end{enumerate}
\end{problem}

\begin{solution}
\begin{enumerate}
\item 由递推关系得 \(a_2=2S_1=2\). 对 \(n\geqslant2\)，有
\[
a_{n+1}-a_n=2S_n-2S_{n-1}=2a_n
\]
即 \(a_{n+1}=3a_n\). 因此
\(a_n=2\cdot3^{n-2}\)，\((n\geqslant2)\).
再结合 \(a_1=1\)，得 \(a_1=1\)，且当 \(n\geqslant2\) 时 \(a_n=2\cdot3^{n-2}\).
\item 当 \(n\geqslant2\) 时，
\[
T_n=\sum_{k=1}^{n}ka_k
 =1+2\sum_{k=2}^{n}k3^{k-2}
\]
记
\[
U_n=\sum_{k=2}^{n}k3^{k-2}
\]
将 \(3U_n\) 的下标向后移一项后与 \(U_n\) 相减，得
\[
2U_n=n3^{n-1}-2-\sum_{j=3}^{n}3^{j-2}
 =n3^{n-1}-2-\frac{3^{n-1}-3}{2}
\]
\[
2U_n=\frac{(2n-1)3^{n-1}-1}{2}
\]
所以
\[
U_n=\sum_{k=2}^{n}k3^{k-2}
 =\frac{(2n-1)3^{n-1}-1}{4}
\]
所以
\[
T_n=\frac{(2n-1)3^{n-1}+1}{2}
\]
当 \(n=1\) 时，上式也给出 \(T_1=1\)，故
\(T_n=\frac{(2n-1)3^{n-1}+1}{2}\)，\((n\in\N^*)\).
\end{enumerate}
\end{solution}

\begin{problem}
如图，已知点 \(F(1, 0)\)，直线 \(l: x = -1\)，\(P\) 为平面上的动点，过 \(P\) 作直线 \(l\) 的垂线，垂足为点 \(Q\)，且 \(\overrightarrow{QP} \cdot \overrightarrow{QF} = \overrightarrow{FP} \cdot \overrightarrow{FQ}\).

\begin{enumerate}
\item 求动点 \(P\) 的轨迹 \(C\) 的方程.
\item 过点 \(F\) 的直线交轨迹 \(C\) 于 \(A\)，\(B\) 两点，交直线 \(l\) 于点 \(M\).

\begin{enumerate}
\item 已知 \(\overrightarrow{MA} = \lambda_1 \overrightarrow{AF}\)，\(\overrightarrow{MB} = \lambda_2 \overrightarrow{BF}\)，求 \(\lambda_1 + \lambda_2\) 的值；
\item 求 \(\left| \overrightarrow{MA} \right| \cdot \left| \overrightarrow{MB} \right|\) 的最小值.
\end{enumerate}
\end{enumerate}

\bitmapfigure[max width=0.28\linewidth]{img/2007/fujian_liberal/q22_fig1.png}
\end{problem}

\begin{solution}
\begin{enumerate}
\item 设 \(P=(x,y)\). 因为 \(PQ\perp l\)，而 \(l:x=-1\)，所以 \(Q=(-1,y)\). 于是
\(\overrightarrow{QP}=(x+1,0)\)，\(\overrightarrow{QF}=(2,-y)\)
\(\overrightarrow{FP}=(x-1,y)\)，\(\overrightarrow{FQ}=(-2,y)\).
代入已知数量积关系，得
\[
2(x+1)=-2(x-1)+y^2
\]
即
\[
y^2=4x
\]
故动点 \(P\) 的轨迹 \(C\) 的方程为 \(y^2=4x\).
\item 设过 \(F\) 的直线斜率为 \(k\). 因为该直线与 \(l\) 相交且与轨迹交于两个不同的点，所以 \(k\neq0\)，其方程可写为
\[
y=k(x-1)
\]
设它与抛物线的交点横坐标为 \(x_1\)，\(x_2\)，则
\[
k^2(x-1)^2=4x
\]
即
\[
k^2x^2-(2k^2+4)x+k^2=0
\]
由韦达定理，\(x_1+x_2=2+\frac{4}{k^2}\)，\(x_1x_2=1\).
因为两交点不同，\(x_1\)，\(x_2\) 是上述方程的两个实根；又 \(x_1x_2>0\) 且 \(x_1+x_2>0\)，所以 \(x_1\)，\(x_2>0\).
\begin{enumerate}
\item 直线与 \(l\) 的交点为 \(M=(-1,-2k)\). 设这两个交点为 \(X_i=(x_i,k(x_i-1))\)（\(i=1\)，\(2\)），并令 \(X_1=A\)，\(X_2=B\). 对于横坐标为 \(x_i\) 的交点，有
\[
\overrightarrow{MX_i}=(x_i+1)(1,k)
\]
\[
\overrightarrow{X_iF}=(1-x_i)(1,k)
\]
由 \(\overrightarrow{MX_i}=\lambda_i\overrightarrow{X_iF}\)，得
\[
\lambda_i=\frac{x_i+1}{1-x_i}=-1+\frac{2}{1-x_i}
\]
因此
\[
\lambda_1+\lambda_2=-2+2\left(\frac{1}{1-x_1}+\frac{1}{1-x_2}\right)
 =-2+2\cdot\frac{2-(x_1+x_2)}{1-(x_1+x_2)+x_1x_2}=0
\]
故 \(\lambda_1+\lambda_2=0\).
\item 因为
\[
|MX_i|=|x_i+1|\sqrt{1+k^2}
\]
所以
\[
|\overrightarrow{MA}|\cdot|\overrightarrow{MB}|
 =(1+k^2)(x_1+1)(x_2+1)
 =4(1+k^2)\left(1+\frac{1}{k^2}\right)
\]
\[
|\overrightarrow{MA}|\cdot|\overrightarrow{MB}|
 =4\left(k^2+2+\frac{1}{k^2}\right)\geqslant16
\]
由基本不等式 \(k^2+\dfrac{1}{k^2}\geqslant2\)，当且仅当 \(k^2=1\) 时取等号，因此 \(\left|\overrightarrow{MA}\right|\cdot\left|\overrightarrow{MB}\right|\) 的最小值为 \(16\).
\end{enumerate}
\end{enumerate}
\end{solution}
